\scanpage{34}
(\mspage{2}

On a alors un $A$ isomorphisme $A'=A''/J$, où $A''=A\inserted{[t_1,}\ldots t_n\inserted{]}$ et $J$ est un idéal de $A''$. Ecrivant $J$ comme limite inductive de ses sous-idéaux $J_i$ de type fini, et posant $A'_i=A''/J_i$, on aura encore $A'=\mathop{\lim}\limits_{\inserted{\longrightarrow}\atop i}A'_i$ donc $B'$ provient d'une algèbre étale $B'_i$ sur un $A'_i$ (loco citato). D'autre part, comme $\operatorname{Spec}A'\to\operatorname{Spec}A$ est surjectif, il en est de même des $\operatorname{Spec}A'_i\to\operatorname{Spec}A$, de plus on vérifie facilement que puisque $\operatorname{Spec}A'\to\operatorname{Spec}A$ est radiciel, il en est de même de $\operatorname{Spec}A'_i\to\operatorname{Spec}A$ pour $i$ assez grand (appliquer EGA IV 1.9.9. à $\operatorname{Spec}A$). Cela nous ramène au cas où $A'$ est un des $A'_i$, donc de présentation finie sur $A$.

Cela prouve la première assertion de 2.1. Le fait que $f_*$ et $f^*$ soient des équivalences quasi-inverses l'une de l'autre en résulte aussitôt.

\uline{Corollaire} 1.2. Soient $F$ un faisceau abélien sur $X$, $F'$ son image inverse sur $X'$, alors les homomorphismes canoniques
\[
(*)\qquad H^i(X,F)\longrightarrow H^i(X',F')
\]
sont des isomorphismes. De même si $F'$ est un faisceau abélien sur $X'$, $F$ son image sur $X$, les homomorphismes canoniques $(*)$ sont des isomorphismes.

\inserted{\uline{Exemples}}
1.3. Voici des exemples d'application fréquents de 1.1. :

\margin{\figure{\begin{tikzpicture}[x=1cm,y=1cm]
\node[draw,ellipse,rotate=23,align=center] (n) at (0,0) {\unsure{laisser blancs}};
\draw (0.9,0.5)--(1.8,0.85)--(2.15,0.85)--(2.15,1.15)--(2.55,1.15);
\draw (1,0)--(1.8,-0.25)--(2.15,-0.25)--(2.15,0.05)--(2.55,0.05);
\draw (0.8,-0.45)--(1.8,-1.5)--(2.15,-1.5)--(2.15,-1.2)--(2.55,-1.2);
\end{tikzpicture}}}

a) $X'$ est un sous-préschéma de $X$ ayant même ensemble sous-jacent que $X$, i.e. défini par un nil-Idéal $I$.

b) $X$ est un préschéma sur un corps séparablement clos $k$, $k'$ est une clôture algébrique de $k$, et $X'=X\otimes_k k'$.

\deleted{\uline{\unsure{Remarques}}} c) Soit $X$ un préschéma géométriquement unibranche (par exemple une courbe algébrique sur un corps $k$, n'ayant comme singularités que des singularités ``cuspidales'', à l'exclusion en particulier de points doubles ordinaires). Alors, si $X'$ est le \unsure{normalis}
\margin{33}
